\section{Amostragem Natural}

\begin{frame}{Amostragem Natural}{Contexto}
	
	A amostragem ideal utiliza um trem de impulsos de Dirac $\delta(t)$, que é uma abstração matemática puramente teórica e fisicamente inexequível, pois exige pulsos com largura infinitesimal ($\tau \to 0$) e amplitude infinita para reter energia finita. 
	
	A amostragem natural surge para tornar o processo fisicamente realizável: em vez de impulsos instantâneos, utiliza-se um trem de pulsos retangulares de largura finita $\tau$ e período $T_s$, gerados por chaves eletrônicas reais (como transistores operando em comutação) que se fecham por um intervalo $\tau$ a cada ciclo. Ela aproxima o sistema de um cenário real de chaveamento analógico.
	
\end{frame}

% Definição do Pulso p0(t)
\begin{frame}{Amostragem Natural}{Definição do Pulso Fundamental $p_0(t)$}
	
	O processo de amostragem natural utiliza pulsos retangulares de largura finita $\tau$. Definimos inicialmente o pulso unitário $p_0(t)$ centrado na origem:
	
	\begin{equation*}
		p_0(t) = 
		\begin{cases}
			1, & -\frac{\tau}{2} \le t \le \frac{\tau}{2} \\[6pt]
			0, & \text{caso contrário}
		\end{cases}
	\end{equation*}
\end{frame}

\begin{frame}{Amostragem Natural}{Definição do Pulso Fundamental $p_0(t)$}
	\vfill
	\centering
	\begin{tikzpicture}
		\begin{axis}[
			pamaxis, 
			ylabel={$p_0(t)$}, 
			xmin=-0.6, xmax=0.6, 
			ymin=-0.2, ymax=1.4,
			xlabel={$t$}
			]
			% Pulso p0(t)
			\draw[sinalcor, line width=1.2pt, fill=sinalcor!20] 
			(axis cs:-0.2, 0) -- (axis cs:-0.2, 1) -- (axis cs:0.2, 1) -- (axis cs:0.2, 0) -- cycle;
			
			% Marcações do eixo t (-tau/2 e tau/2)
			\node[below, font=\small] at (axis cs:-0.2, 0) {$-\frac{\tau}{2}$};
			\node[below, font=\small] at (axis cs:0.2, 0) {$\frac{\tau}{2}$};
			\node[left, font=\small] at (axis cs:0, 1.1) {$1$};
		\end{axis}
	\end{tikzpicture}
\end{frame}

% Definição do Trem de Pulsos p(t)
\begin{frame}{Amostragem Natural}{Trem Periódico de Pulsos $p(t)$}
	
	O sinal de amostragem $p(t)$ é construído pela replicação periódica do pulso $p_0(t)$ a cada período de amostragem $T_s$:
	
	\begin{equation*}
		p(t) = \sum_{k=-\infty}^{\infty} p_0(t - k T_s)
	\end{equation*}
	
\end{frame}

\begin{frame}{Amostragem Natural}{Trem Periódico de Pulsos $p(t)$}
	\vfill
	\centering
	\begin{tikzpicture}
		\begin{axis}[pamaxis, ylabel={$p(t)$}, ymin=-0.2, ymax=1.4]
			\addplot[draw=none] {0};
			\pgfplotsextra{
				\foreach \n in {0,...,\Nam}{
					\pgfmathsetmacro\xa{\n*\Ts - \Tau/2}
					\pgfmathsetmacro\xb{\n*\Ts + \Tau/2}
					\draw[sinalcor, line width=1pt, fill=sinalcor!20] 
					(axis cs:\xa,0) rectangle (axis cs:\xb,1);
				}
			}
			\marcaTs
		\end{axis}
	\end{tikzpicture}
	\par\vspace{0.3em}
	{\small Trem de pulsos periódico com período $T_s$ e largura de pulso $\tau$.}
\end{frame}

% Conexão e Série de Fourier de p(t)
\begin{frame}{Amostragem Natural}{Expansão em Série de Fourier de $p(t)$}
	Como $p(t)$ é estritamente periódico com período $T_s$ ($f_s = 1/T_s$), podemos representá-lo via Série Exponencial de Fourier:
	\begin{equation*}
		p(t) = \sum_{k=-\infty}^{\infty} P[k] e^{j 2\pi k f_s t}
	\end{equation*}
	
	O cálculo dos coeficientes $P[k]$ é realizado no intervalo de um período $[-\frac{T_s}{2}, \frac{T_s}{2}]$:
	\begin{align*}
		P[k] &= \frac{1}{T_s} \int_{-T_s/2}^{T_s/2} p_0(t) e^{-j 2\pi k f_s t} \, dt = \frac{1}{T_s} \int_{-\tau/2}^{\tau/2} (1) e^{-j 2\pi k f_s t} \, dt \\[6pt]
		&= \frac{1}{T_s} \left[ \frac{e^{-j 2\pi k f_s t}}{-j 2\pi k f_s} \right]_{-\tau/2}^{\tau/2} = \frac{1}{T_s} \left( \frac{e^{j \pi k f_s \tau} - e^{-j \pi k f_s \tau}}{j 2\pi k f_s} \right)
	\end{align*}
\end{frame}

\begin{frame}{Amostragem Natural}{Coeficientes $P[k]$ da Série de Fourier}
	Aplicando a identidade de Euler $\sin(\theta) = \frac{e^{j\theta} - e^{-j\theta}}{2j}$:
	\begin{equation*}
		P[k] = \frac{\sin(\pi k f_s \tau)}{\pi k T_s f_s}
	\end{equation*}
	
	Como $T_s f_s = 1$, multiplicamos e dividimos por $\tau$ para obter a função $\text{sinc}$:
	\begin{equation*}
		P[k] = \frac{\tau}{T_s} \cdot \frac{\sin(\pi k f_s \tau)}{\pi k f_s \tau} \implies P[k] = \frac{\tau}{T_s} \text{sinc}(k f_s \tau)
	\end{equation*}
	
	\begin{block}{Sinal de Amostragem no Tempo}
		\begin{equation*}
			p(t) = \sum_{k=-\infty}^{\infty} \left[ \frac{\tau}{T_s} \text{sinc}(k f_s \tau) \right] e^{j 2\pi k f_s t}
		\end{equation*}
	\end{block}
\end{frame}

% Construção Analítica de s(t) no Tempo
\begin{frame}{Amostragem Natural}{Modelagem Analítica no Tempo}
	O sinal amostrado $s(t)$ é gerado pelo produto direto entre $x(t)$ e $p(t)$:
	
	\begin{equation*}
		s(t) = x(t) \cdot p(t)
	\end{equation*}
	
	Substituindo a expansão em Série de Fourier de $p(t)$:
	
	\begin{equation*}
		s(t) = x(t) \cdot \sum_{k=-\infty}^{\infty} P[k] e^{j 2\pi k f_s t}
	\end{equation*}
	
	
\end{frame}

\begin{frame}{Amostragem Natural}{Modelagem Analítica no Tempo}
	Pela propriedade distributiva, inserimos $x(t)$ no somatório:
	
	\begin{equation*}
		s(t) = \sum_{k=-\infty}^{\infty} P[k] x(t) e^{j 2\pi k f_s t}
	\end{equation*}
	
	
	Substituindo a expressão analítica de $P[k]$:
	
	\begin{equation*}
		s(t) = \sum_{k=-\infty}^{\infty} \left[ \frac{\tau}{T_s} \text{sinc}(k f_s \tau) \right] x(t) e^{j 2\pi k f_s t}
	\end{equation*}
\end{frame}

% Amostragem Natural: sobreposição x(t) + p(t), e produto s(t) = x(t)·p(t)
\begin{frame}{Amostragem Natural}{Sobreposição e Produto}
	\centering
	\begin{minipage}[c][0.72\textheight][c]{0.95\textwidth}
		\centering
		
		% Parâmetros locais (Ts = 0.2 s, Tau = 0.05 s, 5 pulsos)
		\def\Ts{0.2}
		\def\Tau{0.05}
		\def\Nam{5}
		
		\only<1>{
			\begin{tikzpicture}
				\begin{axis}[
					pamaxis,
					ylabel={$x(t)$ e $p(t)$},
					ylabel style={at={(axis cs:0,1.45)}, anchor=south}, % Centraliza sobre o eixo x=0
					trig format=rad,
					ymin=-0.2, ymax=1.45
					]
					% Trem de pulsos p(t)
					\pgfplotsextra{%
						\foreach \n in {0,...,\Nam}{%
							\pgfmathsetmacro\xa{\n*\Ts - \Tau/2}%
							\pgfmathsetmacro\xb{\n*\Ts + \Tau/2}%
							\draw[gray!60, line width=.7pt, fill=gray!15]
							(axis cs:\xa,0) rectangle (axis cs:\xb,1);
						}%
					}
					% x(t) = cos(2*pi*2*t)
					\addplot[sinalcor, line width=1.4pt, smooth, domain=0:1, samples=300]
					{cos(4*pi*x)};
				\end{axis}
			\end{tikzpicture}\par
			\small Sinal analógico $x(t)=\cos(2\pi\cdot 2t)$ sobreposto ao trem de pulsos $p(t)$ (largura $\tau$, período $T_s$).}
		
		% produto s(t)
		\only<2>{%
			\begin{tikzpicture}
				\begin{axis}[
					pamaxis,
					ylabel={$s(t) = x(t) \cdot p(t)$},
					ylabel style={at={(axis cs:0,1.45)}, anchor=south}, % Centraliza sobre o eixo x=0
					trig format=rad,
					ymin=-0.2, ymax=1.45
					]
					% Envoltória x(t) como referência tracejada
					\addplot[gray!60, dashed, line width=.7pt, smooth, domain=0:1, samples=300]
					{cos(4*pi*x)};
					% Pulsos de amostragem natural: topo SEGUE x(t)
					\pgfplotsextra{%
						\foreach \n in {0,...,\Nam}{%
							\pgfmathsetmacro\xa{\n*\Ts - \Tau/2}%
							\pgfmathsetmacro\xb{\n*\Ts + \Tau/2}%
							\draw[sinalcor, line width=.9pt, fill=sinalcor!25]
							(axis cs:\xa,0) --
							(axis cs:\xa,{cos(4*pi*\xa)}) --
							plot[domain=\xa:\xb, samples=25, smooth]
							(axis cs:\x,{cos(4*pi*\x)}) --
							(axis cs:\xb,{cos(4*pi*\xb)}) --
							(axis cs:\xb,0) -- cycle;
						}%
					}
				\end{axis}
			\end{tikzpicture}\par
			\small Sinal amostrado naturalmente $s(t)=x(t)\cdot p(t)$. O topo de cada pulso \textbf{acompanha a forma de $x(t)$} — característica da amostragem natural.}
		
		% ---------------------------------------------------- (3) destaque: onde p(t)=0 vs p(t)=1
		\only<3>{%
			\begin{tikzpicture}
				\begin{axis}[
					pamaxis,
					ylabel={$s(t)$},
					ylabel style={at={(axis cs:0,1.45)}, anchor=south}, % Centraliza sobre o eixo x=0
					trig format=rad,
					ymin=-0.2, ymax=1.45
					]
					% x(t) como referência tracejada
					\addplot[gray!60, dashed, line width=.7pt, smooth, domain=0:1, samples=300]
					{cos(4*pi*x)};
					% Regiões onde p(t) = 0  ->  s(t) = 0 (faixas cinzas claras)
					\pgfplotsextra{%
						\foreach \n in {0,...,\Nam}{%
							\pgfmathsetmacro\xa{\n*\Ts + \Tau/2}%
							\pgfmathsetmacro\xb{(\n+1)*\Ts - \Tau/2}%
							\ifnum\n<\Nam
							\draw[black!10, fill=black!5]
							(axis cs:\xa,-0.15) rectangle (axis cs:\xb,1.35);
							\fi
						}%
					}
					% Pulsos de amostragem natural
					\pgfplotsextra{%
						\foreach \n in {0,...,\Nam}{%
							\pgfmathsetmacro\xa{\n*\Ts - \Tau/2}%
							\pgfmathsetmacro\xb{\n*\Ts + \Tau/2}%
							\draw[sinalcor, line width=.9pt, fill=sinalcor!25]
							(axis cs:\xa,0) --
							(axis cs:\xa,{cos(4*pi*\xa)}) --
							plot[domain=\xa:\xb, samples=25, smooth]
							(axis cs:\x,{cos(4*pi*\x)}) --
							(axis cs:\xb,{cos(4*pi*\xb)}) --
							(axis cs:\xb,0) -- cycle;
						}%
					}
					% Rótulos das duas regiões
					\node[font=\tiny, anchor=north] at (axis cs:0.30,-0.02)
					{$p(t)=0 \Rightarrow s(t)=0$};
					\node[font=\tiny, anchor=north, text=sinalcor] at (axis cs:0.05,-0.02)
					{$p(t)=1$};
				\end{axis}
			\end{tikzpicture}\par
			\small Onde $p(t)=0$, o produto é nulo (faixas cinzas). Onde $p(t)=1$ (dentro de cada pulso), o produto reproduz $x(t)$ — o topo \textbf{segue} o cosseno.}
		
	\end{minipage}
\end{frame}

% Domínio da Frequência: Convolução e Espectro Final
\begin{frame}{Amostragem Natural}{Análise no Domínio da Frequência}
	
	A multiplicação no tempo equivale à convolução no domínio da frequência:
	
	\begin{equation*}
		S(f) = \mathcal{F}\{s(t)\} = X(f) * P(f)
	\end{equation*}
	
	A Transformada de Fourier do trem de pulsos $p(t)$ é:
	
	\begin{equation*}
		P(f) = \sum_{k=-\infty}^{\infty} P[k] \delta(f - k f_s)
	\end{equation*}
	
\end{frame}

\begin{frame}{Amostragem Natural}{Análise no Domínio da Frequência}
	Calculando a convolução $X(f) * P(f)$:
	\begin{align*}
		S(f) &= X(f) * \sum_{k=-\infty}^{\infty} P[k] \delta(f - k f_s) \\[4pt]
		&= \sum_{k=-\infty}^{\infty} P[k] \left[ X(f) * \delta(f - k f_s) \right] \\[4pt]
		&= \sum_{k=-\infty}^{\infty} P[k] X(f - k f_s)
	\end{align*}
\end{frame}

% Frame 2: Espectro Final do Sinal Amostrado
\begin{frame}{Amostragem Natural}{Espectro Final do Sinal Amostrado}
	
	Substituindo o coeficiente $P[k] = \frac{\tau}{T_s} \text{sinc}(k f_s \tau)$, temos a expressão espectral final:
	
	\begin{equation*}
		S(f) = \sum_{k=-\infty}^{\infty} \frac{\tau}{T_s} \text{sinc}(k f_s \tau) X(f - k f_s)
	\end{equation*}
	
	\vfill
	\centering
	\begin{tikzpicture}
		\begin{axis}[
			pamaxis, 
			ylabel={$S(f)$}, 
			xlabel={$f$},
			xmin=-15, xmax=15, 
			ymin=-0.1, ymax=0.6,
			axis lines=middle,
			trig format=rad
			]
			% Envelope Sinc
			% Réplicas Espectrais centradas em k*fs
		\end{axis}
	\end{tikzpicture}
	\par\vspace{0.3em}
	{\small Réplicas de $X(f)$ repetidas a cada $f_s$, atenuadas pelo envelope $\frac{\tau}{T_s}\text{sinc}(k f_s \tau)$.}
\end{frame}